\chapter*{Preface}
\addcontentsline{toc}{chapter}{Preface}
\markboth{Preface}{Preface}

This volume presents the complete visual classification, combinatorial taxonomy, and machine-readable relational inventory of composition mosaics across the circular rings $\mathbb{Z}_1$ through $\mathbb{Z}_9$.

As the third monograph in our coffee-table trilogy on the combinatorial geometry of circular lattices---standing alongside \emph{The Polygonal Burnsidean: An Atlas of Inscribed Polygons from $\mathbb{Z}_1$ to $\mathbb{Z}_{16}$} and \emph{Hexadecality: The Atlas of Hexadecagonal Compositions}---this atlas unifies the full progression of cyclic lattices from the elementary monad to the rich enneagonal universe:
\begin{itemize}
    \item Across the first nine rings $\mathbb{Z}_1$ through $\mathbb{Z}_9$, the combinatorial universe generates exactly $26{,}442$ set partitions, governed by the initial segment of the Bell numbers $B_1 = 1, B_2 = 2, B_3 = 5, B_4 = 15, B_5 = 52, B_6 = 203, B_7 = 877, B_8 = 4{,}140$, and $B_9 = 21{,}147$.
    \item Under the physical symmetries of the respective dihedral groups $D_n$, these $26{,}442$ partitions collapse into exactly $\mathbf{1{,}859}$ non-equivalent dihedral composition mosaics, realising the initial terms of sequence OEIS A084708:
    \begin{equation*}
    1, \; 2, \; 3, \; 7, \; 12, \; 37, \; 93, \; 354, \; 1{,}350 \quad \implies \quad \sum_{n=1}^9 M_n = 1{,}859
    \end{equation*}
    \item The nine lattices capture the full structural spectrum of cyclic geometry: the trivial monad $\mathbb{Z}_1$; the reflexive dyad $\mathbb{Z}_2$; the odd prime rings $\mathbb{Z}_3, \mathbb{Z}_5, \mathbb{Z}_7$; the composite square lattice $\mathbb{Z}_4$; the composite hexagonal lattice $\mathbb{Z}_6$ where scalene triangles and chiral enantiomorphic blocks first emerge into the geometric universe; the cubic power-of-two lattice $\mathbb{Z}_8$, which hosts the unique cyclic Skolem difference sequence of order $4$; and the nonagonal lattice $\mathbb{Z}_9$, whose ternary modular geometry exhibits rich threefold rotational symmetries and modular prime-power factorisations.
    \item Through the two-vertex domino inflation functor, the 8-vertex Skolem partition acts as the seed for higher-order affine macro-orbits on the 16-cycle $\mathbb{Z}_{16}$, creating an explicit mathematical bridge connecting all three volumes of the trilogy.
\end{itemize}

\section*{Cartographic Conventions and Visual Aesthetics}

Each mosaic in this atlas is presented in pure geometric form. The partition blocks appear as shaded convex inscribed polygons, metric chords, and vertex pips within the circular frame of the regular $n$-gon, unencumbered by peripheral coordinate clutter:
\begin{itemize}
    \item \textbf{Achiral Blocks (Blue)}: Inscribed polygons possessing an internal axis of bilateral reflection symmetry across the cyclic lattice.
    \item \textbf{Chiral Blocks (Pink)}: Scalene polygons that lack reflectional symmetry, existing as one half of an enantiomorphic pair.
\end{itemize}

The complete visual catalogue in Chapter~\ref{chap:visual-atlas} displays all $1{,}859$ configurations arranged in a strict taxonomic sequence across 33 high-density vector plates. All accompanying analytical invariants---dihedral orbit sizes $|\mathcal{O}|$, stabiliser subgroups, Fa\`a di Bruno partition multiplicities, chord dispersion energies, polygon profiles, and relational manifests linking directly into \emph{The Polygonal Burnsidean}---are systematically catalogued in Chapter~\ref{chap:canonical-taxonomy} and in the companion machine-readable open datasets.

\bigskip
\begin{flushright}
\textbf{Gemma Knee} \qquad \textbf{Antigra Vitti} \\
\emph{Newcastle upon Tyne} \\
\emph{November 2026}
\end{flushright}
